\chapter{Finding the Rhythm}

\nopagebreak
\epigraph{"I've never had a schedule for this blog: I write when I feel like writing, which usually means several times a week."}{"Blogging pace" (2024)}
\nopagebreak

Temporal patterns are autobiography written in timestamps. Every post carries a date and time. Aggregate 5,233 of them and a portrait appears: when one person does his thinking.

\bigskip

\section*{The Tuesday Peak}

Tuesday leads most years, though Wednesday takes the top spot in several.

\begin{longtable}{p{3.00cm}p{3.00cm}p{3.00cm}}
\toprule
\textbf{Day} & \textbf{Posts} & \textbf{\% of Total} \tabularnewline
\midrule
Tuesday & 1,003 & 19.2\% \tabularnewline
Wednesday & 938 & 17.9\% \tabularnewline
Monday & 887 & 17.0\% \tabularnewline
Thursday & 823 & 15.7\% \tabularnewline
Friday & 561 & 10.7\% \tabularnewline
Saturday & 555 & 10.6\% \tabularnewline
Sunday & 466 & 8.9\% \tabularnewline
\bottomrule
\end{longtable}

The pattern is stable. Tuesday and Wednesday compete for the top spot across years, with Tuesday winning overall. The weekend always trails. Not a hobby blogger working when inspiration strikes on lazy Sundays. A professional integrating writing into a workweek.

Why Tuesday? Perhaps Monday is for catching up on email and meetings. By Tuesday, the week has settled. There's time to write.

The Tuesday pattern survived career changes, platform changes, and seventeen years of life. It's an invariant: something about how Cook works that transcends context.

\bookfigure{ch03_weekday_radial.png}{Weekly rhythm}

\bigskip

\section*{The 7 AM Hour}

The peak hour is 7:00 to 7:59 AM.

\begin{longtable}{p{3.00cm}p{3.00cm}p{3.00cm}}
\toprule
\textbf{Hour} & \textbf{Posts} & \textbf{\% of Total} \tabularnewline
\midrule
7 AM & 683 & 13.1\% \tabularnewline
8 AM & 580 & 11.1\% \tabularnewline
6 AM & 501 & 9.6\% \tabularnewline
9 AM & 503 & 9.6\% \tabularnewline
\bottomrule
\end{longtable}

Forty-three percent of all posts appear between 6 AM and 10 AM. The blog is a morning practice.

The pattern suggests writing before the day's demands take over. Before meetings, before email, before the work that doesn't leave room for reflection. The 7 AM post is the first intellectual act of the day.

Compare the morning cluster (6-10 AM: 2,267 posts) to the afternoon cluster (12-5 PM: 1,226 posts) to the evening cluster (6-10 PM: 770 posts). Morning wins by nearly two to one.

\bigskip

\section*{The Habit Before the Content}

This rhythm existed from the beginning. The 7 AM peak is visible in 2008. The Tuesday lead is visible in 2008. The weekend slump is visible in 2008.

The habit preceded the content. Before Cook knew what the blog would become, he knew when he would write it.

\begin{blogquote}
"The hard part isn't having ideas. The hard part is sitting down at the same time every day."
--- (Editorial observation)
\end{blogquote}

This may be the secret to seventeen years. The rhythm was established early and protected fiercely. The content could change, from PowerShell to Python, from clinical trials to cryptography, but the rhythm stayed constant. Tuesday mornings. 7 AM. Before the rest of the day.

\bigskip

\section*{Monthly Flatness}

Unlike days and hours, months show no pattern.

\begin{longtable}{p{3.00cm}p{3.00cm}p{3.00cm}}
\toprule
\textbf{Month} & \textbf{Posts} & \textbf{\%} \tabularnewline
\midrule
January & 469 & 9.0\% \tabularnewline
February & 444 & 8.5\% \tabularnewline
March & 436 & 8.3\% \tabularnewline
April & 406 & 7.8\% \tabularnewline
May & 394 & 7.5\% \tabularnewline
June & 413 & 7.9\% \tabularnewline
July & 410 & 7.8\% \tabularnewline
August & 438 & 8.4\% \tabularnewline
September & 448 & 8.6\% \tabularnewline
October & 480 & 9.2\% \tabularnewline
November & 478 & 9.1\% \tabularnewline
December & 417 & 8.0\% \tabularnewline
\bottomrule
\end{longtable}

Each month accounts for 7.5-9.2\% of posts. No summer slump. No December hiatus. No April surge.

This flatness is the signature of consistency. Seasonal variation would suggest a practice that competes with other life rhythms: vacations, holidays, academic calendars. The flatness suggests a practice that has become routine, independent of external calendars.

\bigskip

\section*{The Gaps}

Every blogger takes breaks. The question is how long.

The longest gap in the entire archive is 11 days: July 9 to July 21, 2014.

The second-longest gaps are 8 days, appearing three times.

No gap reaches two weeks. No sabbatical. No "I'm back after six months." In seventeen years, the blog has never gone silent for more than eleven days.

Most professionals take longer vacations than the blog's longest gap. A two-week holiday would set a new record. A month away would be unprecedented. Whatever else was happening in Cook's life---travel, illness, family---the blog kept appearing.

\bigskip

\section*{What Rhythm Reveals}

The temporal patterns reveal a practice designed to survive.

\textbf{Morning writing} gives the blog the best hours---before fatigue and distraction.

\textbf{Weekday dominance} integrates the practice with professional life rather than competing with it.

\textbf{Monthly flatness} means the blog doesn't depend on mood or season.

\textbf{Short gaps} prevent restart friction from accumulating.

Each property makes the practice more sustainable. Morning writing prevents falling behind. Weekday integration avoids the choice between work and blog. Monthly consistency avoids rebuilding momentum.

The rhythm is engineered for persistence. It's not the rhythm of inspiration. It's the rhythm of showing up.

\bigskip

\section*{The Practice as Foundation}

The habit established in 2008-2012 was not just about writing. It was about when to write.

The rhythm that emerged in the first year---Tuesday mornings, 7 AM, consistent across months---became the foundation for everything that followed. When the content changed, the rhythm remained. When the voice evolved, the rhythm remained. When interests shifted, the rhythm remained.

This is what sustained practice looks like underneath: not just what you do, but when you do it, protected across years.

\bigskip

\section*{So What?}

\textit{The rhythm was established before the content clarified. The sustainability came from engineering the practice around life's constraints—mornings before work, weekdays that fit professional rhythms, gaps short enough to prevent restart friction—not from waiting for inspiration. Seventeen years later, the content had transformed completely, but the Tuesday morning habit remained.}

